\scanpage{9}
qu'on ne peut avoir simultanément
\[
\{x,y,z\},\ \{x,y,-z\}\in\widetilde F_M
\]
en effet cela implique $x+y+z=x+y-z=0$ d'où $2z=0$, d'où (comme $z\in M^*$) que $2\,1_\Lambda=0$, mais alors $\varepsilon\unsure{=|\ne}\varepsilon'=\varepsilon''$ quand-même contrairement à l'hypothèse.

On peut donc identifier $F_M$ à une partie de $\unsure{\mathfrak P|\mathscr P}_3(S_M)$. Si on suppose \uline{$\mathrm{Pic}(\Lambda)=0$}, on trouve encore
% slanted
\margin{\unsure{inutile?|inutile}, tout élément de $M^*$ fait tjrs partie d'une base}
\[
S_M=\bigcup_{A\in\mathfrak P_3(S_M)}A\tag*{(19)}
\]
% slanted
\margin{Si on n'a pas $\mathrm{Pic}(\Lambda)=0$, alors il vaut mieux définir $S_M=\widetilde S'_M/\{\pm1\}$, où $\widetilde S'_M\subset\widetilde S_M$ est \unsure{déf.|défini} comme $\bigcup_{A\in\widetilde F_M}A$.}

Je désigne par $A_M$ l'ens. des parties à deux éléments \inserted{de $S_M$,} contenues dans un $\deleted{A}f\in F_M$ convenable
\[
\begin{aligned}
A_M&=\{a\in\mathfrak P_2(S_M)\mid\exists f\in F_M,\ a\subset f\}.\\
&=\left\{a\in\mathfrak P_2(S_M)\ \middle|\ 
\begin{gathered}\exists(x,y)\in\mathrm{Rep}(M),\ \text{t.q.}\\ \{\dot x,\dot y\}=a\end{gathered}\right\}
\end{aligned}\tag*{(20)}
\]
\deleted{\unsure{f}} \deleted{Soit}
% slanted
\margin{On \unsure{suppose|supposera} $2.1_\Lambda\ne0$, sinon la situation est \unsure{purement} triviale $\widetilde S_M=S_M,\ \widetilde F_M=F_M$, et $f$ est unique.}

soit $a=\{\dot x,\dot y\}\in A_M$, avec $(x,y)\in\mathrm{Rep}(M)$, je cherche à déterminer l'ens. des $f\in F_M$ qui contiennent $a$. On doit donc avoir $f=\dot A$, où $A$ est de la forme
% archive sheet number: 8
